% Written by Anders Sjoqvist and Ulf Lundstrom, 2009
% The main sources are: tinyKACTL, Beta and Wikipedia

\section{数学}

\subsection{方程}
\[ax^2+bx+c=0 \Rightarrow x = \frac{-b\pm\sqrt{b^2-4ac}}{2a}\]

极值点为 $x = -b/2a$.

\[\begin{aligned}ax+by=e\\cx+dy=f\end{aligned}
\Rightarrow
\begin{aligned}x=\dfrac{ed-bf}{ad-bc}\\y=\dfrac{af-ec}{ad-bc}\end{aligned}\]

一般地, 对方程 $Ax = b$, 变量 $x_i$ 的解为
\[x_i = \frac{\det A_i'}{\det A} \]
其中 $A_i'$ 是把 $A$ 的第 $i$ 列换成 $b$ 所得的矩阵.

\subsection{递推}
若 $a_n = c_1 a_{n-1} + \dots + c_k a_{n-k}$, 且 $r_1, \dots, r_k$ 是 $x^k - c_1 x^{k-1} - \dots - c_k$ 的互异根, 则存在 $d_1, \dots, d_k$ 使得
\[a_n = d_1r_1^n + \dots + d_kr_k^n. \]
重根 $r$ 会乘上多项式因子, 如 $a_n = (d_1n + d_2)r^n$.

\subsection{三角函数}
\begin{align*}
\sin(v+w)&{}=\sin v\cos w+\cos v\sin w\\
\cos(v+w)&{}=\cos v\cos w-\sin v\sin w\\
\end{align*}
\begin{align*}
\tan(v+w)&{}=\dfrac{\tan v+\tan w}{1-\tan v\tan w}\\
\sin v+\sin w&{}=2\sin\dfrac{v+w}{2}\cos\dfrac{v-w}{2}\\
\cos v+\cos w&{}=2\cos\dfrac{v+w}{2}\cos\dfrac{v-w}{2}
\end{align*}
\[ (V+W)\tan(v-w)/2{}=(V-W)\tan(v+w)/2 \]
其中 $V, W$ 是角 $v, w$ 的对边长.
\begin{align*}
	a\cos x+b\sin x&=r\cos(x-\phi)\\
	a\sin x+b\cos x&=r\sin(x+\phi)
\end{align*}
其中 $r=\sqrt{a^2+b^2}, \phi=\operatorname{atan2}(b,a)$.

\subsection{几何}

\subsubsection{三角形}
边长: $a,b,c$\\
半周长: $p=\dfrac{a+b+c}{2}$\\
面积: $A=\sqrt{p(p-a)(p-b)(p-c)}$\\
外接圆半径: $R=\dfrac{abc}{4A}$\\
内切圆半径: $r=\dfrac{A}{p}$\\
中线长 (把三角形分成面积相等的两半): $m_a=\tfrac{1}{2}\sqrt{2b^2+2c^2-a^2}$\\
角平分线长: $s_a=\sqrt{bc\left[1-\left(\dfrac{a}{b+c}\right)^2\right]}$\\
正弦定理: $\dfrac{\sin\alpha}{a}=\dfrac{\sin\beta}{b}=\dfrac{\sin\gamma}{c}=\dfrac{1}{2R}$\\
余弦定理: $a^2=b^2+c^2-2bc\cos\alpha$\\
正切定理: $\dfrac{a+b}{a-b}=\dfrac{\tan\dfrac{\alpha+\beta}{2}}{\tan\dfrac{\alpha-\beta}{2}}$\\

\subsubsection{四边形}
设边长 $a,b,c,d$, 对角线 $e, f$, 对角线夹角 $\theta$, 面积 $A$,
记 $F=b^2+d^2-a^2-c^2$:

\[ 4A = 2ef \cdot \sin\theta = F\tan\theta = \sqrt{4e^2f^2-F^2} \]

 圆内接四边形对角之和为 $180^\circ$,
$ef = ac + bd$, 且 $A = \sqrt{(p-a)(p-b)(p-c)(p-d)}$.

\subsubsection{球坐标}
\begin{center}
\includegraphics[width=25mm]{content/math/sphericalCoordinates}
\end{center}
\[\begin{array}{cc}
x = r\sin\theta\cos\phi & r = \sqrt{x^2+y^2+z^2}\\
y = r\sin\theta\sin\phi & \theta = \textrm{acos}(z/\sqrt{x^2+y^2+z^2})\\
z = r\cos\theta & \phi = \textrm{atan2}(y,x)
\end{array}\]

\subsection{导数/积分}
\begin{align*}
	\dfrac{d}{dx}\arcsin x = \dfrac{1}{\sqrt{1-x^2}} &&& \dfrac{d}{dx}\arccos x = -\dfrac{1}{\sqrt{1-x^2}} \\
	\dfrac{d}{dx}\tan x = 1+\tan^2 x &&& \dfrac{d}{dx}\arctan x = \dfrac{1}{1+x^2} \\
	\int\tan ax = -\dfrac{\ln|\cos ax|}{a} &&& \int x\sin ax = \dfrac{\sin ax-ax \cos ax}{a^2} \\
	\int e^{-x^2} = \frac{\sqrt \pi}{2} \text{erf}(x) &&& \int xe^{ax}dx = \frac{e^{ax}}{a^2}(ax-1)
\end{align*}

分部积分:
\[\int_a^bf(x)g(x)dx = [F(x)g(x)]_a^b-\int_a^bF(x)g'(x)dx\]

\subsection{求和}
\[ c^a + c^{a+1} + \dots + c^{b} = \frac{c^{b+1} - c^a}{c-1}, c \neq 1 \]
\begin{align*}
	1 + 2 + 3 + \dots + n &= \frac{n(n+1)}{2} \\
	1^2 + 2^2 + 3^2 + \dots + n^2 &= \frac{n(2n+1)(n+1)}{6} \\
	1^3 + 2^3 + 3^3 + \dots + n^3 &= \frac{n^2(n+1)^2}{4} \\
	1^4 + 2^4 + 3^4 + \dots + n^4 &= \frac{n(n+1)(2n+1)(3n^2 + 3n - 1)}{30} \\
\end{align*}

\subsection{级数}
$$e^x = 1+x+\frac{x^2}{2!}+\frac{x^3}{3!}+\dots,\,(-\infty<x<\infty)$$
$$\ln(1+x) = x-\frac{x^2}{2}+\frac{x^3}{3}-\frac{x^4}{4}+\dots,\,(-1<x\leq1)$$
$$\sqrt{1+x} = 1+\frac{x}{2}-\frac{x^2}{8}+\frac{2x^3}{32}-\frac{5x^4}{128}+\dots,\,(-1\leq x\leq1)$$
$$\sin x = x-\frac{x^3}{3!}+\frac{x^5}{5!}-\frac{x^7}{7!}+\dots,\,(-\infty<x<\infty)$$
$$\cos x = 1-\frac{x^2}{2!}+\frac{x^4}{4!}-\frac{x^6}{6!}+\dots,\,(-\infty<x<\infty)$$

\subsection{概率论}
设 $X$ 是离散随机变量, 取值 $x$ 的概率为 $p_X(x)$. 其期望 (均值) 为 $\mu=\mathbb{E}(X)=\sum_xxp_X(x)$, 方差为 $\sigma^2=V(X)=\mathbb{E}(X^2)-(\mathbb{E}(X))^2=\sum_x(x-\mathbb{E}(X))^2p_X(x)$, 其中 $\sigma$ 是标准差. 若 $X$ 是连续的, 则有概率密度函数 $f_X(x)$, 上面的求和改为积分, 并把 $p_X(x)$ 换成 $f_X(x)$.

期望是线性的:
\[\mathbb{E}(aX+bY) = a\mathbb{E}(X)+b\mathbb{E}(Y)\]
若 $X$ 与 $Y$ 独立, \[V(aX+bY) = a^2V(X)+b^2V(Y).\]

\subsection{Markov 链}
\emph{Markov 链} 是一种离散随机过程, 下一状态只依赖当前状态.
设 $X_1,X_2,\ldots$ 是该过程产生的随机变量序列.
则有转移矩阵 $\mathbf{P} = (p_{ij})$, 其中 $p_{ij} = \Pr(X_n = i | X_{n-1} = j)$,
且 $\mathbf{p}^{(n)} = \mathbf P^n \mathbf p^{(0)}$ 是 $X_n$ 的概率分布 (即 $p^{(n)}_i = \Pr(X_n = i)$),
其中 $\mathbf{p}^{(0)}$ 是初始分布.

\paragraph{平稳分布}
若 $\mathbf{\pi} = \mathbf{\pi P}$, 则称 $\mathbf{\pi}$ 为平稳分布.
若 Markov 链\emph{不可约} (任意状态都能到达任意状态),
则 $\pi_i = \frac{1}{\mathbb{E}(T_i)}$, 其中 $\mathbb{E}(T_i)$ 是相邻两次访问状态 $i$ 的期望间隔.
$\pi_j/\pi_i$ 是相邻两次访问 $i$ 之间访问状态 $j$ 的期望次数.

在连通, 无向, 非二分的图上, 若等概率走向各邻居, 则 $\pi_i$ 与点 $i$ 的度数成正比.

\paragraph{遍历性}
若极限分布与初始分布无关, 则称 Markov 链是\emph{遍历的}.
有限 Markov 链遍历当且仅当它不可约且\emph{非周期} (即各环长的 gcd 为 1).
$\lim_{k\rightarrow\infty}\mathbf{P}^k = \mathbf{1}\pi$.

\paragraph{吸收}
若状态可分成两个集合 $\mathbf{A}$ 和 $\mathbf{G}$, 使 $\mathbf{A}$ 中的状态都是吸收态 ($p_{ii}=1$), 且 $\mathbf{G}$ 中的状态都能到达 $\mathbf{A}$ 中的吸收态, 则称其为 A-链.
从 $j$ 出发, 被状态 $i\in\mathbf{A}$ 吸收的概率为 $a_{ij} = p_{ij}+\sum_{k\in\mathbf{G}} a_{ik}p_{kj}$.
从 $i$ 出发, 到被吸收的期望时间为 $t_i = 1+\sum_{k\in\mathbf{G}}p_{ki}t_k$.

\subsection{期望与概率技巧}
\begin{compactitem}
\item \textbf{期望线性性}: 不要求独立. 求个数的期望就拆成指示变量, $\mathbb{E}[\text{个数}]=\sum_i \Pr[\text{事件 } i]$.
\item \textbf{尾和}: $X$ 取非负整数时 $\mathbb{E}[X]=\sum_{k\ge1}\Pr[X\ge k]$. 独立同分布的最大值: $\Pr[\max\le k]=F(k)^n$.
\item \textbf{期望 DP 倒推}: $E_i=c_i+\sum_j p_{ij}E_j$, 终点 $E=0$, 按拓扑序从后往前算.
  只有自环时移项: $E_i=(c_i+\sum_{j\ne i}p_{ij}E_j)/(1-p_{ii})$.
\item \textbf{有环, 高斯消元}: 把 $E_i-\sum_j p_{ij}E_j=c_i$ 列成方程组, 用 \texttt{SolveLinear}, $O(n^3)$, $n$ 约到 500.
  取模时把除法换成乘逆元. 吸收概率同理.
\item \textbf{只有一个未知量}: 设起点期望为 $x$, 其他状态都写成 $a_i+b_ix$ 一路推下去, 最后解 $x=a+bx$, 不用高斯消元.
\item \textbf{二分答案}: 可以随时重来的题, 答案 $x$ 出现在自己的转移里, 即 $E_i=\min(\text{继续}, x)$.
  对猜的 $x$ 做一遍 DP 得到 $f(x)$; 若 $f(x)<x$ 则真答案 $<x$, 否则 $\ge x$. 二分约 100 次.
\item \textbf{方差}: DP 里同时维护 $\mathbb{E}[X]$ 和 $\mathbb{E}[X^2]$, 如 $\mathbb{E}[(X+1)^2]=\mathbb{E}[X^2]+2\mathbb{E}[X]+1$.
\item \textbf{min-max 容斥}: $\mathbb{E}[\max S]=\sum_{\emptyset\ne T\subseteq S}(-1)^{|T|+1}\mathbb{E}[\min T]$. 用于全部出现的期望时间.
\item \textbf{常用结论}: 成功率 $p$ 的试验, 期望 $1/p$ 次成功一次. 集齐 $n$ 种卡片: $nH_n$.
  $n$ 个 $U[0,1]$ 中第 $k$ 小的期望: $k/(n+1)$.
  随机排列: 逆序对期望 $n(n-1)/4$, 不动点期望 $1$, 环数期望 $H_n$.
\item \textbf{随机游走}: 树上从 $u$ 走到父亲的期望步数是 $2\cdot\mathit{size}(u)-1$.
  $0$ 处反射, 从 $0$ 走到 $n$ 期望 $n^2$ 步. 公平赌徒从 $i$ 出发: 先到 $n$ 的概率 $i/n$, 期望步数 $i(n-i)$.
\item \textbf{精度}: 期望很大时用 \texttt{long double}; 方程太大时可以反复迭代 $E\leftarrow c+PE$ 直到收敛.
\end{compactitem}

